\documentclass[11pt,a4paper]{article}
\usepackage[utf8]{inputenc}
\usepackage[T1]{fontenc}
\usepackage[spanish]{babel}
\usepackage[margin=2.2cm]{geometry}
\usepackage{amsmath,amssymb}
\usepackage{enumitem}
\usepackage{xcolor}
\usepackage{titlesec}
\usepackage{array,booktabs,longtable}
\usepackage{listings}
\usepackage{hyperref}

\hypersetup{colorlinks=true,urlcolor=blue,linkcolor=black}
\setlist{itemsep=2pt,topsep=2pt}
\titlespacing*{\section}{0pt}{10pt}{4pt}
\titlespacing*{\subsection}{0pt}{8pt}{3pt}

\definecolor{codebg}{rgb}{0.96,0.96,0.96}
\lstset{
  basicstyle=\ttfamily\footnotesize,
  backgroundcolor=\color{codebg},
  frame=single,
  framesep=4pt,
  breaklines=true,
  columns=fullflexible,
  showstringspaces=false,
  literate={ñ}{{\~n}}1 {á}{{\'a}}1 {é}{{\'e}}1 {í}{{\'i}}1 {ó}{{\'o}}1 {ú}{{\'u}}1
}

\title{\vspace{-1cm}\textbf{TP5 -- Sistema 1: Implementación del motor de simulación}\\[2pt]\large Kuramoto -- Java}
\author{Simulación de Sistemas -- 2026Q1}
\date{}

\begin{document}
\maketitle
\vspace{-0.8cm}

Motor de simulación del modelo de Kuramoto en Java, modularizado en \textbf{6 archivos} dentro de \texttt{motor/} (total $\sim 7.6$ kB, bajo el límite de 20 kB del enunciado). Sólo motor: el análisis y la animación se hacen aparte sobre los CSV generados.

\section{Estructura del código}

Seis clases, una por archivo, cada una con una sola responsabilidad. Sin dependencias externas (solo \texttt{java.io} y \texttt{java.util}).

\begin{center}
\renewcommand{\arraystretch}{1.2}
\begin{tabular}{@{}p{3.4cm} p{2.6cm} p{8.6cm}@{}}
\toprule
\textbf{Archivo} & \textbf{Clase} & \textbf{Responsabilidad} \\
\midrule
\texttt{Config.java} & \texttt{Config} (record) & Lleva todos los parámetros + parseo de CLI (\texttt{fromArgs}). \\
\texttt{Network.java} & \texttt{Network} + enum \texttt{Topology} & Construye la lista de vecinos (\texttt{int[][]}) según topología. \\
\texttt{Dynamics.java} & \texttt{Dynamics} & Evalúa $f(\theta)_i = \omega_i + K\sum_{j\in\text{vec}(i)}\sin(\theta_j-\theta_i)$. \\
\texttt{Integrator.java} & \texttt{Integrator} & Un paso de RK4. Buffers $\mathbf k_1..\mathbf k_4$, \texttt{tmp} preasignados. \\
\texttt{Output.java} & \texttt{Output} & Cabecera + filas del CSV (computa $r(t)$ en cada fila). \\
\texttt{KuramotoSim.java} & \texttt{KuramotoSim} & \texttt{main}: orquesta inicialización, loop y volcado. \\
\bottomrule
\end{tabular}
\end{center}

\subsection*{Dependencias entre módulos (sin ciclos)}

\begin{lstlisting}
KuramotoSim ---+--> Config
               +--> Network ---> Topology
               +--> Integrator ---> Dynamics
               +--> Output ---> Config
\end{lstlisting}

\subsection*{Estructura de datos: listas de vecinos}

En vez de guardar $A_{ij}$ como matriz densa $N\times N$ (250\,000 entradas para $N=500$), guardamos para cada nodo $i$ un \textbf{arreglo de índices de sus vecinos}. Esto: (i) recorre sólo los vecinos reales en \texttt{deriv} (rápido para redes esparsas), (ii) es equivalente para el caso completo, (iii) aprovecha que el grafo no cambia durante la simulación.

\subsection*{Reproducibilidad: dos semillas}

\begin{itemize}
\item \texttt{--seed}: controla $\omega_i$ y $\theta_i(0)$. Misma semilla $\Rightarrow$ mismas condiciones iniciales \emph{independientemente de la topología}.
\item \texttt{--netSeed}: controla el sorteo de $A_{ij}$ (sólo \texttt{random}). Default = \texttt{seed}.
\end{itemize}

Facilita comparar topologías sobre la misma condición inicial y promediar separadamente sobre realizaciones de red y de condiciones iniciales.

\section{Parámetros (CLI)}

Todos parametrizables vía \texttt{--clave valor}.

\subsection*{Parámetros del modelo}

\begin{center}
\renewcommand{\arraystretch}{1.2}
\begin{tabular}{@{}l l l p{7.7cm}@{}}
\toprule
\textbf{Flag} & \textbf{Tipo} & \textbf{Default} & \textbf{Descripción} \\
\midrule
\texttt{--N} & int & 500 & Número de osciladores (consigna: $N>500$). \\
\texttt{--K} & double & 1.0 & Acoplamiento global $K\in[0,1]$. \\
\texttt{--topology} & string & complete & \texttt{complete}/\texttt{random}/\texttt{ring}. \\
\texttt{--p} & double & 0.5 & Probabilidad de conexión (sólo \texttt{random}). \\
\texttt{--v} & int & 1 & Vecindad del anillo (sólo \texttt{ring}), $v\in[1,10]$. \\
\texttt{--muOmega} & double & 1.0 & Media de la normal para $\omega_i$. \\
\texttt{--sigmaOmega} & double & 0.1 & Desvío estándar de la normal para $\omega_i$. \\
\bottomrule
\end{tabular}
\end{center}

\subsection*{Parámetros numéricos y de salida}

\begin{center}
\renewcommand{\arraystretch}{1.2}
\begin{tabular}{@{}l l l p{7.7cm}@{}}
\toprule
\textbf{Flag} & \textbf{Tipo} & \textbf{Default} & \textbf{Descripción} \\
\midrule
\texttt{--dt} & double & 0.01 & Paso de integración (fijo). Probar $10^{-x}$, $x\in\{1,\ldots,4\}$. \\
\texttt{--tSim} & double & 100.0 & Tiempo total simulado. \\
\texttt{--seed} & long & 42 & Semilla para $\omega_i$ y $\theta_i(0)$. \\
\texttt{--netSeed} & long & = seed & Semilla para $A_{ij}$ (sólo \texttt{random}). \\
\texttt{--output} & string & kuramoto.csv & Archivo de salida. \\
\texttt{--dumpEvery} & int & 1 & Vuelca cada $n$ pasos al CSV. \\
\texttt{--dumpPhases} & bool & true & Si \texttt{false}, sólo \texttt{t,r} (mucho más liviano). \\
\bottomrule
\end{tabular}
\end{center}

\section{Formato del archivo de salida (CSV)}

\begin{lstlisting}
# N=500 K=1.000000 dt=0.010000 tSim=100.000000 topology=COMPLETE p=0.500000 v=1 muOmega=1.000000 sigmaOmega=0.100000 seed=42 netSeed=42 dumpEvery=1
t,r,theta_0,theta_1,...,theta_{N-1}
0.0,0.0763...,1.4234...,3.1415...,...
0.01,0.0791...,1.4334...,3.1525...,...
...
\end{lstlisting}

\begin{itemize}
\item Primera línea: comentario con todos los parámetros (autodocumenta el archivo).
\item Segunda línea: nombres de columna.
\item Filas siguientes: valores. Si \texttt{--dumpPhases false}, sólo \texttt{t} y \texttt{r}.
\end{itemize}

$r(t)$ se computa dentro del motor en \texttt{writeRow} ($O(N)$ por fila): el postproceso lo lee directo, y \texttt{--dumpPhases false} habilita un modo ``solo $r$'' liviano para los barridos masivos.

\section{Compilación y ejecución}

\textbf{Configuración del TP} (calibrada experimentalmente, ver \texttt{calibracion.pdf}):
$N=600$, $dt=10^{-3}$, $t_{\text{sim}}=50$ (completa, aleatoria) / $1500$ (anillo).

\begin{lstlisting}[language=bash]
cd motor
javac *.java   # compila las 6 clases en .class

# Red completa, K=0.5
java KuramotoSim --N 600 --K 0.5 --topology complete \
                 --dt 0.001 --tSim 50 \
                 --seed 42 --output complete_K0.5_seed42.csv \
                 --dumpEvery 10 --dumpPhases true

# Red aleatoria, K=0.1, p=0.1, semillas separadas
java KuramotoSim --N 600 --K 0.1 --topology random --p 0.1 \
                 --dt 0.001 --tSim 50 \
                 --seed 42 --netSeed 7 \
                 --output random_p0.1_K0.1_s42_n7.csv \
                 --dumpPhases false

# Red anillo, v=3  -- tSim mucho mayor por la fisica del anillo
java KuramotoSim --N 600 --K 0.7 --topology ring --v 3 \
                 --dt 0.001 --tSim 1500 \
                 --seed 42 --output ring_v3_K0.7.csv
\end{lstlisting}

\section{Algoritmo}

\textbf{Inicialización:}
\begin{enumerate}
\item $\omega_i \sim \mathcal{N}(1, 0.1^2)$ con la RNG de \texttt{seed} (vía \texttt{nextGaussian()}).
\item $\theta_i(0) \sim \mathcal{U}[0, 2\pi)$ con la misma RNG.
\item Lista de vecinos según topología (con RNG de \texttt{netSeed} si corresponde).
\end{enumerate}

\textbf{Loop principal -- RK4:} con $f(\theta)_i = \omega_i + K \sum_{j\in\text{vec}(i)} \sin(\theta_j-\theta_i)$,

\begin{align*}
\mathbf k_1 &= f(\boldsymbol\theta), \quad \mathbf k_2 = f(\boldsymbol\theta + \tfrac{dt}{2}\mathbf k_1), \quad \mathbf k_3 = f(\boldsymbol\theta + \tfrac{dt}{2}\mathbf k_2), \quad \mathbf k_4 = f(\boldsymbol\theta + dt\,\mathbf k_3) \\
\boldsymbol\theta &\leftarrow \boldsymbol\theta + \tfrac{dt}{6}\bigl(\mathbf k_1 + 2\mathbf k_2 + 2\mathbf k_3 + \mathbf k_4\bigr)
\end{align*}

El sistema es autónomo; \texttt{deriv} no recibe argumento de tiempo. Los buffers $\mathbf k_1..\mathbf k_4$ y \texttt{tmp} se alocan una sola vez fuera del loop.

\textbf{Cálculo de $r$:}
\[
c_x = \frac{1}{N}\sum_j \cos\theta_j, \quad c_y = \frac{1}{N}\sum_j \sin\theta_j, \quad r = \sqrt{c_x^2 + c_y^2}
\]

\textbf{Construcción de vecinos:}
\begin{itemize}
\item \emph{complete:} \texttt{nbr[i] = \{0..N-1\} \textbackslash\ \{i\}}.
\item \emph{random:} para cada par $(i,j)$ con $i<j$, sortear \texttt{nextDouble() < p}; si sí, agregar $j$ a \texttt{nbr[i]} y $i$ a \texttt{nbr[j]} (red \textbf{no dirigida} por convención: el enunciado no lo exige explícitamente).
\item \emph{ring:} \texttt{nbr[i] = \{(i-v) mod N, ..., (i+v) mod N\} \textbackslash\ \{i\}} (exactamente $2v$ vecinos). Aritmética modular con \texttt{((i-d) \% N + N) \% N} para manejar negativos en Java.
\end{itemize}

\section{Lo que el motor \emph{no} hace (a propósito)}

\begin{itemize}
\item No grafica ni anima.
\item No calcula $\tau_s$; emite $r(t)$, el umbral y la detección del estacionario van en postproceso.
\item No promedia sobre realizaciones (ejecutar el motor varias veces con distintas \texttt{--seed}).
\item No barre $K$, $p$, $v$ automáticamente (cada combinación se corre como proceso separado, fácil de paralelizar con \texttt{xargs -P} o un script).
\end{itemize}

Esto sigue el requisito del enunciado: \emph{``el análisis y módulo de animación se ejecuta en forma independiente tomando estos archivos de texto como input.''}

\end{document}
